
Si @@M0@@,
alors @@M1@@ est surjectif par 1.3 (ii); donc @@M2@@ est surjectif. Par conséquent, @@M3@@. D'autre part, si @@M4@@ n'est pas surjectif, alors @@M5@@; donc @@M6@@ est \uline{strictement} décroissant depuis @@M7@@, jusqu'au moment de devenir @@M8@@; d'où (iii).

\textbf{Définition 1.5.} \emph{Soient @@M9@@ un faisceau cohérent sur @@M10@@, @@M11@@ un entier @@M12@@ et @@M13@@ une suite de @@M14@@ entiers. Alors @@M15@@ est dit un \emph{@@M16@@-faisceau} s'il satisfait à un des deux ensembles équivalents suivants de conditions}~:

@@M17@@
\begin{enumerate}[label=(\roman*)]
\item \emph{En chaque point du support de @@M18@@, @@M19@@ est engendré par @@M20@@.}
\item \emph{@@M21@@.}
\item \emph{(Lorsque @@M22@@), il existe une section @@M23@@ qui définit une suite exacte @@M24@@ telle que @@M25@@ soit un @@M26@@-faisceau.}
\end{enumerate}

@@M27@@ \emph{Il existe une suite @@M28@@-régulière de @@M29@@ sections @@M30@@ telle que, pour @@M31@@, si @@M32@@ désigne la restriction de @@M33@@ au schéma des zéros de @@M34@@, alors @@M35@@.}

\textbf{Proposition 1.6.} \emph{Soit @@M36@@ un @@M37@@-faisceau sur @@M38@@. Alors}~:
\begin{enumerate}[label=(\roman*)]
\item \emph{Pour toute suite ``assez générale'' @@M39@@ de @@M40@@ sections @@M41@@, si @@M42@@ désigne la restriction de @@M43@@ au schéma des zéros de @@M44@@, alors @@M45@@, où @@M46@@.}
\item \emph{Tout sous-faisceau @@M47@@ de @@M48@@ est un @@M49@@-faisceau.}
\end{enumerate}

